% v2.7.0 figure UID: FIG-P715-01
\tikzset{slfig-FIG-P715-01/.style={font=\fontsize{9.5pt}{11.5pt}\selectfont,every path/.append style={line cap=round,line join=round}}}
\pgfkeysifdefined{/tikz/alt}{}{\tikzset{alt/.code={}}}
\begin{figure}[htbp]
\centering
\begin{tikzpicture}[slfig-FIG-P715-01,
  >=Stealth,
  every node/.style={font=\fontsize{9.5pt}{11.5pt}\selectfont,align=center},
  panel/.style={draw=SLRuleGray,line width=0.56pt,rounded corners=3pt},
  title/.style={font=\fontsize{10.4pt}{12.4pt}\selectfont\bfseries,text=SLDeepBlue,inner sep=0pt},
  page/.style={circle,draw=SLDeepBlue,fill=SLSoftBlue,line width=0.72pt,minimum size=9mm,inner sep=1pt,font=\fontsize{10.2pt}{12.2pt}\selectfont\bfseries},
  link/.style={->,draw=SLMidBlue,line width=0.72pt},
  chosen/.style={->,draw=SLAccentOrange,line width=1.10pt},
  edge note/.style={font=\fontsize{9.5pt}{11.5pt}\selectfont,inner sep=0pt},
  formula/.style={font=\fontsize{12pt}{14pt}\selectfont,inner sep=0pt},
  cell/.style={draw=SLRuleGray,line width=0.48pt,minimum width=6.6mm,minimum height=5.8mm,inner sep=0pt,font=\fontsize{10.2pt}{12.2pt}\selectfont},
  note/.style={font=\fontsize{9.5pt}{11.5pt}\selectfont,text=SLTextGray,inner sep=0pt},
  focus/.style={draw=SLAccentOrange,line width=1.00pt,inner sep=0pt},
  >={Stealth[length=2.15mm,width=1.55mm]},
  alt={对象—关系—结论：按结点顺序(i,j,h)的单位权网页图含i→j、j→i、j→h、h→i，出度为(1,2,1)；邻接矩阵A按来源列存边，M按出度归一且列和为1，P=M的转置且行和为1，两种状态推进由概率向量转置相连。}]

\useasboundingbox (-7.90,-3.80) rectangle (7.90,3.30);
\draw[panel] (-7.90,-3.75) rectangle (0.35,3.25);
\draw[panel] (0.65,-3.75) rectangle (7.90,3.25);
\node[title] at (-3.78,2.91) {网页图、邻接矩阵与列归一};
\node[title] at (4.28,2.91) {行随机转置桥};

\node[page] (gi) at (-6.60,1.50) {$i$};
\node[page] (gj) at (-4.10,1.50) {$j$};
\node[page] (gh) at (-5.35,0.15) {$h$};
\draw[chosen] (gj) to[bend right=15]
  node[edge note,above=1.2mm] {$j\to i$} (gi);
\draw[link] (gi) to[bend right=15]
  node[edge note,below=1.2mm] {$i\to j$} (gj);
\draw[link] (gj) -- (gh);
\draw[link] (gh) -- (gi);

\node[note] at (-1.70,1.75) {四条边权均为 $1$};
\node[note] at (-1.70,1.28) {矩阵行、列顺序均为 $(i,j,h)$};
\node[formula] at (-1.70,0.78) {$A_{ij}>0\Longleftrightarrow j\to i$};
\node[formula] at (-1.70,0.12)
  {$c_j=\sum_r A_{rj}$\\[-1pt]$(c_i,c_j,c_h)=(1,2,1)$};

\node[formula,anchor=east] at (-6.60,-1.50) {$A=$};
\matrix (ma) [matrix of nodes,nodes={cell},
  row sep=-\pgflinewidth,column sep=-\pgflinewidth]
  at (-5.35,-1.50)
  {$0$&$1$&$1$\\$1$&$0$&$0$\\$0$&$1$&$0$\\};
\node[focus,fit=(ma-1-2)] {};

\node[formula,anchor=east] at (-2.72,-1.50) {$M=$};
\matrix (mm) [matrix of nodes,nodes={cell},
  row sep=-\pgflinewidth,column sep=-\pgflinewidth]
  at (-1.48,-1.50)
  {$0$&$1/2$&$1$\\$1$&$0$&$0$\\$0$&$1/2$&$0$\\};
\node[focus,fit=(mm-1-2)] {};
\node[formula] at (-5.35,-2.72) {$M_{:j}=A_{:j}/c_j$};
\node[formula] at (-1.48,-2.72)
  {$\boldsymbol1^{\mathsf T}M=\boldsymbol1^{\mathsf T}$};
\node[formula] at (-3.78,-3.30)
  {$\boldsymbol p^{(t+1)}=M\boldsymbol p^{(t)}$};

\node[note] at (4.28,2.34) {同一结点顺序 $(i,j,h)$};
\node[formula,anchor=east] at (2.92,1.35) {$P=$};
\matrix (mp) [matrix of nodes,nodes={cell},
  row sep=-\pgflinewidth,column sep=-\pgflinewidth]
  at (4.28,1.35)
  {$0$&$1$&$0$\\$1/2$&$0$&$1/2$\\$1$&$0$&$0$\\};
\node[focus,fit=(mp-2-1)] {};
\node[formula] at (4.28,0.18) {$P=M^{\mathsf T}$};
\node[formula] at (4.28,-0.38) {$P_{ji}=M_{ij}$};
\node[formula] at (4.28,-0.94)
  {$=\Pr(X_{t+1}=i\mid X_t=j)$};
\node[formula] at (4.28,-1.55) {$P\boldsymbol1=\boldsymbol1$};
\node[formula] at (4.28,-2.15) {$\rho_{t+1}=\rho_tP$};
\node[formula] at (4.28,-2.76)
  {$\rho_t=(\boldsymbol p^{(t)})^{\mathsf T}$};
\end{tikzpicture}
\caption{列随机约定下的网页有向图，并给出与行随机马尔可夫链约定的显式转置桥。}
\label{fig:V5-C07-web-random-walk}
\end{figure}
